% from ega4-16-fr.tex
\oldpage[IV-4]{55}

\begin{remarks}[16.11.3]
\EGAReaderTarget{ega4.c04.statement.16-11-3}{IV.16.11.3}
\begin{enumerate}
  \item[(i)] Le fait que les $D_{\mathbf{p}}$ sont deux à deux permutables en
  vertu de (16.11.2.2) n'implique naturellement pas que la
  $\mathcal{O}_U$-Algèbre $\operatorname{Diff}_{U/S}$ soit commutative, les
  $D_{\mathbf{p}}$ ne permutant aux produits par les sections de
  $\mathcal{O}_U$ que si $n=0$.
  \item[(ii)] Les indices $\mathbf{p}$ tels que $|\mathbf{p}|=1$ sont les
  $\boldsymbol{\varepsilon}_\lambda
  =(\varepsilon_{\lambda\mu})_{\mu\in L}$ avec
  $\varepsilon_{\lambda\mu}=0$ si $\mu\neq\lambda$ et
  $\varepsilon_{\lambda\lambda}=1$; lorsque $L$ est fini, les opérateurs
  $D_{\boldsymbol{\varepsilon}_\lambda}$ ne sont autres que les
  $S$-dérivations $D_i$ introduites dans (16.5.7). On notera qu'en général
  (et contrairement à ce qui se passe en Analyse classique), il n'est pas
  vrai qu'un opérateur différentiel d'ordre quelconque puisse s'écrire comme
  combinaison linéaire de puissances des $D_i$ (cf. (16.12)).
  \item[(iii)] Pour tout entier $r\geq1$, on peut définir la notion de
  morphisme \emph{différentiellement lisse jusqu'à l'ordre $r$} en remplaçant
  dans (16.10.1) la condition (ii) par la condition que les homomorphismes
  \[
    \mathbf{S}_{\mathcal{O}_X}^m(\Omega_{X/S}^1)
    \longrightarrow\mathscr{G}r_m(\mathscr{P}_{X/S})
  \]
  soient bijectifs \emph{pour tout $m\leq r$}. Le raisonnement de (16.11.2)
  prouve alors que si, dans la condition a), on remplace «différentiellement
  lisse» par «différentiellement lisse jusqu'à l'ordre $r$», cette condition
  est équivalente à la condition b) dans laquelle on se borne aux
  $\mathbf{p}\in\mathbb{N}^{(L)}$, $\mathbf{q}\in\mathbb{N}^{(L)}$ tels que
  $|\mathbf{p}|\leq r$, $|\mathbf{q}|\leq r$.
\end{enumerate}
\end{remarks}

\subsection{Cas de la caractéristique nulle: critère jacobien pour les
morphismes différentiellement lisses}
\EGAReaderTarget{ega4.c04.subsec.16-12}{IV.16.12}

\begin{env}[16.12.1]
\EGAReaderTarget{ega4.c04.statement.16-12-1}{IV.16.12.1}
On dit qu'un préschéma $X$ est \emph{de caractéristique $p$} ($p$ égal à $0$
ou à un nombre premier) si, pour tout ouvert affine $U$ de $X$, l'anneau
$\Gamma(U,\mathcal{O}_X)$ est de caractéristique $p$ (0, 21.1.1). Il résulte
de (0, 21.1.3) que pour que $X$ soit de caractéristique $0$, il faut et il
suffit que pour tout point \emph{fermé} $x$ de $X$, le corps résiduel $k(x)$
soit de caractéristique $0$, ou encore que $X$ puisse être muni d'une
structure de $\mathbb{Q}$-préschéma (nécessairement unique).
\end{env}

\begin{theorem}[16.12.2]
\label{IV.16.12.2}
Soient $X$ un préschéma de caractéristique $0$, $f:X\to S$ un morphisme. Si
$\Omega_{X/S}^1$ est un $\mathcal{O}_X$-Module localement libre (non
nécessairement de type fini), $f$ est différentiellement lisse.
\end{theorem}

La question étant locale sur $X$, on peut supposer qu'il existe une famille
$(z_\lambda)$ de sections de $\mathcal{O}_X$ au-dessus de $X$ telle que
$(dz_\lambda)$ soit une base du $\mathcal{O}_X$-Module $\Omega_{X/S}^1$.
Appliquant le critère (16.11.2), il suffit de remarquer que les opérateurs
\[
  D_{\mathbf{p}}=(\mathbf{p}!)^{-1}
  \prod_\lambda D_\lambda^{p_\lambda}
\]
(où les $D_\lambda$ sont les formes coordonnées correspondant à la base
$(dz_\lambda)$) vérifient les relations (16.11.2.1), ce qui est une
conséquence du fait que les $D_\lambda$ sont des dérivations.

\begin{env}[16.12.3]
\label{IV.16.12.3}
Le théorème précédent n'est plus exact lorsqu'on abandonne l'hypothèse que
$X$ est de caractéristique $0$. Par exemple, si
$S=\operatorname{Spec}(k)$, où $k$ est un corps de caractéristique $p>0$,
$X=\operatorname{Spec}(K)$ où $K=k(\alpha)$ avec
$\alpha\notin k$, $\alpha^p\in k$, on vérifie aussitôt
